\chapter{Methodology}
\label{chap:method}

\section{Experimental protocol}
All settings of the study (data source, competing methods, metrics, number of repetitions and
random seed) are declared in the topic profile \texttt{\ThesisProfileId}. Each configuration is
evaluated over independent repeated runs, and every run is recorded in a result bundle stamped
with the source-code commit and the SHA-256 digest of the input data.

\ifAutoHasRsm
\section{Response surface methodology}
The process factors are explored with a central composite design comprising factorial, axial and
replicated centre points. The response is modelled by the second-order polynomial
\begin{equation}
  \hat{y} = \beta_{0} + \sum_{i} \beta_{i} x_{i} + \sum_{i} \beta_{ii} x_{i}^{2}
            + \sum_{i<j} \beta_{ij} x_{i} x_{j},
\end{equation}
fitted by ordinary least squares. Model significance is assessed by an $F$-test, adequacy by the
lack-of-fit test against pure error, and the stationary point is classified from the eigenvalues
of the quadratic form.
\fi

\ifAutoHasAblation
\section{Ablation design}
The contribution of each core module is isolated by evaluating cumulative or leave-one-out
configurations of the full system under identical conditions.
\fi

\section{Statistical analysis}
For every metric, normality is examined with the Shapiro--Wilk test and homogeneity of variance
with Levene's test. Two-group comparisons are routed to the paired or Student $t$-test, Welch's
$t$-test, the Wilcoxon signed-rank test or the Mann--Whitney $U$ test according to these
diagnostics, and are accompanied by Cohen's $d$ or the rank-biserial correlation. Comparisons of
several groups use one-way ANOVA with Tukey's HSD or the Kruskal--Wallis $H$ test. Significance
markers follow the conventional thresholds stated beneath each table.
